\chapter{Introduction Générale}

\section{Contexte}

L’entrée à l’université représente un enjeu important pour beaucoup d’étudiants, et les concours d’admission sont souvent une étape décisive. Pouvoir prédire les chances de réussite permet non seulement d’aider les candidats à mieux se préparer, mais aussi d’accompagner les établissements dans leur processus de sélection.

\section{Problématique}

Prédire si un candidat va réussir son concours d’entrée est un vrai défi : les données sont nombreuses, variées, et pas toujours très propres. Difficile, dans ces conditions, de bâtir un modèle fiable. Il fallait donc trouver une manière efficace de traiter ces informations pour en tirer des prédictions utiles.

\section{Objectifs}

L’objectif de ce projet était de développer un modèle capable d’estimer les probabilités de réussite au concours, en se basant sur des données académiques comme les scores au GRE ou au TOEFL, la moyenne générale, et d’autres critères pertinents.

\section{Démarche suivie}

Pour y parvenir, on a commencé par rassembler plusieurs jeux de données, puis on en a choisi un — le plus complet et fiable. Ensuite, on a nettoyé et préparé les données, avant de les explorer pour bien comprendre ce qu’elles contenaient. Plusieurs algorithmes de machine learning ont été testés et comparés pour sélectionner celui qui donnait les meilleurs résultats.

\section{Structure du rapport}

Ce mémoire est organisé en quatre parties : 
\begin{itemize}
    \item La première partie présente la collecte et le prétraitement des données.
    \item La deuxième détaille la modélisation et l’évaluation des différents algorithmes.
    \item La troisième analyse les résultats et les performances du modèle retenu.
    \item Enfin, la dernière partie discute des limitations et propose des pistes d’amélioration pour aller plus loin.
\end{itemize}